% ============================================================
% preamble.tex --- настройки презентации (minted edition)
% ============================================================

% --- Шрифты и язык (pdfLaTeX) ---
\usepackage[T2A]{fontenc}
\usepackage[utf8]{inputenc}
\usepackage[russian,english]{babel}

% --- Дополнительные пакеты ---
\usepackage{graphicx}
\usepackage{booktabs}
\usepackage{xcolor}
\usepackage{mwe}       % example-image-a/b для примеров
\usepackage{qrcode}    % QR-код на заключительном слайде

% --- Тема оформления ---
\usetheme{Boadilla}
\setbeamertemplate{frame numbering}[fraction]

% --- Подсветка кода (minted) ---
\usepackage{minted}

% Настройка цветов для фона (по желанию)
\definecolor{codebg}{HTML}{F8F8F8}
\definecolor{codegray}{HTML}{808080}

% Глобальные настройки minted
\setminted{
  bgcolor=codebg,
  fontsize=\footnotesize,
  frame=single,
  framesep=3mm,
  rulecolor=\color{codegray},
  breaklines=true,
  linenos=false,
  tabsize=2
}

% --- Удобная команда для inline-кода ---
\newcommand{\cmd}[1]{\texttt{\textbackslash #1}}

% --- Заметки докладчика ---
% Раскомментируйте ОДНУ из строк, чтобы увидеть \note в PDF:
% \setbeameroption{show notes}
% \setbeameroption{show notes on second screen=right}

% --- Метаданные ---
\title{Презентации в \LaTeX{}}
\subtitle{Методологии и технологии научно-исследовательской деятельности}
\author{Лихач Александр}
\institute{Кафедра ПМИ}
\date{Аспирантский семинар, \today}